% "Beyond the Model" — A0 portrait conference poster.
%
% Build (from the repo root, which every path below is relative to):
%   sbatch rerun/job_m_poster.sh               figures + numbers + typeset
%   sbatch rerun/job_m_poster.sh --tex-only    typeset only
%
% Every data value is a macro from paper_figures/poster/numbers.tex, written by
% Experiments/render_poster_figures.py from the same artifacts and cross-checks as
% the talk. Do not type numbers into this file; the one exception is the published
% Table III row, quoted from the paper.
%
% Placeholders, safe to leave empty:
%   \RepoURL   set it to the GitHub URL to print a QR code in the footer
%   logos      drop PDFs into Experiments/poster/logos/ (guelph, stjoes, mcmaster)
\documentclass{article}
\usepackage[paperwidth=841mm, paperheight=1189mm, margin=0mm]{geometry}
\usepackage{fontspec}
\usepackage{xcolor}
\usepackage{graphicx}
\usepackage{tikz}
\usetikzlibrary{calc}
\usepackage[most]{tcolorbox}
\usepackage{tabularx}
\usepackage{array}
\usepackage{qrcode}
\usepackage{microtype}

\def\RepoURL{}   % e.g. \def\RepoURL{https://github.com/<org>/<repo>}  (\def, not \newcommand: the footer tests it with \ifx)

% ── type ──────────────────────────────────────────────────────────────────────
\setmainfont{SourceSansPro}[Extension=.otf, UprightFont=*-Regular,
  BoldFont=*-Semibold, ItalicFont=*-RegularIt, BoldItalicFont=*-SemiboldIt]
\newfontfamily\serif{SourceSerifPro}[Extension=.otf, UprightFont=*-Semibold,
  BoldFont=*-Bold]
\setmonofont{FiraMono}[Extension=.otf, UprightFont=*-Regular, BoldFont=*-Medium]
\newcommand\body{\fontsize{25}{30}\selectfont}
\renewcommand\small{\fontsize{21}{25}\selectfont}
\newcommand\num[1]{{\ttfamily\bfseries #1}}
% p = value in the body face. Math mode would switch to Computer Modern.
\newcommand\p[1]{\textit{p}\,=\,#1}

% ── colour: the deck's palette ────────────────────────────────────────────────
\definecolor{ink}{HTML}{0E1618}
\definecolor{bodytext}{HTML}{2C3C40}
\definecolor{muted}{HTML}{61767A}
\definecolor{rule}{HTML}{CDD8D9}
\definecolor{ground}{HTML}{EEF2F2}
\definecolor{sunk}{HTML}{E4EAEA}
\definecolor{teal}{HTML}{0D6B6B}
\definecolor{rust}{HTML}{A4552B}
\definecolor{slate}{HTML}{4A6572}
\definecolor{green}{HTML}{3D7A52}
\definecolor{icu}{HTML}{2F4B7C}
\definecolor{mort}{HTML}{9C6B98}
\pagecolor{ground}
\color{bodytext}

\input{paper_figures/poster/numbers.tex}
\newcommand\fig[1]{paper_figures/poster/figures/#1.pdf}

% ── furniture ─────────────────────────────────────────────────────────────────
\setlength\parindent{0pt}
\setlength\parskip{0pt}
\raggedright
\hyphenpenalty=10000
\exhyphenpenalty=10000

% A white panel with a 3 mm accent rule on top, like the deck's panel + rule.
% Inner width is 237 mm in a column; the render script draws figures to that width.
\newtcolorbox{panel}[1][teal]{enhanced, colback=white, frame hidden, arc=0pt,
  boxrule=0pt, boxsep=0pt, left=10mm, right=10mm, top=9mm, bottom=8mm,
  borderline north={3mm}{0pt}{#1}, before skip=0pt, after skip=0pt,
  before upper={\body\raggedright}}

% Numbered heading that states the finding.
\newcommand\heading[3]{%
  {\ttfamily\bfseries\fontsize{32}{38}\selectfont\color{#2}#1}\hspace{6mm}%
  {\serif\fontsize{36}{42}\selectfont\color{ink}#3}\par\vspace{6mm}}

\renewcommand\caption[1]{\par\vspace{3mm}{\small\color{muted}\itshape #1}\par}
\newcommand\gap{\par\vspace{5mm}}
\newcommand\level[2]{\leavevmode{\ttfamily\bfseries\color{#1}#2}}
\newcommand\bul[1]{\par\hangindent=9mm\hangafter=1\makebox[9mm][l]{\color{teal}\textbullet}#1}

% Columns are built twice: once at natural height to report fit, once at the
% column's fixed height with \colgap as \vfill, so all three end level.
\newlength\colheight
\newsavebox\measure
\newcommand\colgap{\par\vspace{10mm}}
% Measured outside the tikzpicture: inside one, TikZ typesets with \nullfont and
% the measurement comes out short.
\newcommand\measurecolumn[3]{% name, width (mm), content macro
  \setbox\measure\vbox{\hsize=#2mm\linewidth=#2mm\renewcommand\colgap{\par\vspace{10mm}}#3}%
  \typeout{POSTER column #1: natural height
    \the\dimexpr\ht\measure+\dp\measure\relax\space of \the\colheight}}
\newcommand\placecolumn[3]{% x (mm), width (mm), content macro
  \node[anchor=north west, inner sep=0pt] at ($(current page.north west)+(#1mm,-\coltop)$)
    {\begin{minipage}[t][\colheight][t]{#2mm}\renewcommand\colgap{\par\vfill}#3\end{minipage}};%
}

% ── layout (mm) ───────────────────────────────────────────────────────────────
% margin 20 · header 20–126 · hook 136–218 · columns 230–977 · band 989–1117
% · footer 1127–1169. Columns 257 wide at x = 20, 292, 564; band boxes 529 + 257.
\newcommand\coltop{230mm}
\setlength\colheight{747mm}

% ─────────────────────────────────────────────────────────────────────────────
% Column 1 — background and method
% ─────────────────────────────────────────────────────────────────────────────
\newcommand\ColumnOne{%
\begin{panel}[ink]
\heading{1}{ink}{Filtering is a modelling choice}
Clinical ML papers open with a filtering pipeline and rarely report what it
removed. Any rule that edits or discards values changes the distribution, so
the question is not \emph{did filtering help?} but \textbf{is it still the same
problem?}
\gap
\bul{It can make a benchmark easier without making a model better.}
\bul{It is not task-neutral: one filter can help one task and hurt another.}
\bul{Pipelines standardise it; almost none characterise what it did.}
\end{panel}
\colgap
\begin{panel}[ink]
\heading{2}{ink}{One cohort, one model, 13 arms}
\includegraphics{\fig{poster_pipeline}}
\gap
Two tasks on the same 168 features: {\color{icu}\textbf{prolonged ICU stay}}
(>\,3 days; \PrevICURaw\ positive) and {\color{mort}\textbf{in-hospital
mortality}} (\PrevMortRaw). Each filter is applied alone to the raw data, and
every arm inherits the same folds, so arms are paired patient by patient.
\end{panel}
\colgap
\begin{panel}[ink]
\heading{3}{ink}{Eleven filters, three levels}
\includegraphics{\fig{poster_levels}}
\gap
{\small
\renewcommand\arraystretch{1.08}
\begin{tabularx}{\linewidth}{@{}>{\raggedright\arraybackslash}p{40mm}>{\raggedright\arraybackslash}X@{}}
\level{teal}{VALUE}   & 7 per-vital ranges, e.g.\ heart rate \ThrHR, temperature
                        \ThrTemp; and their union, \emph{all vitals} \\
\level{slate}{FEATURE} & \emph{fill missing data}: fills empty cells (\ThrFill) \\
\level{rust}{RECORD}  & \emph{long missing segment}: \ThrLMS \\
\level{rust}{RECORD}  & \emph{long gap}: \ThrLG \\
\level{rust}{RECORD}  & \emph{high invalid data}: \ThrHID \\
\end{tabularx}}
\gap
\textbf{Only a record-level filter can delete a patient.}
\end{panel}
\colgap
\begin{panel}[ink]
\heading{4}{ink}{The raw data really is broken}
The heart-rate column peaks at \num{\color{rust}9,999,999}\,bpm, a charting
sentinel rather than a patient. Removing \HRRemoved\ of 1.16~million readings
(\HRPctRemoved) takes the maximum to 459 and the standard deviation from 9,291
to 24.
\gap
\includegraphics{\fig{poster_readings}}
\caption{Implausible readings removed per vital by the value-level filters.}
\gap
Filtering is necessary. The question is what else it does.
\end{panel}
}

% ─────────────────────────────────────────────────────────────────────────────
% Column 2 — what the paper measured
% ─────────────────────────────────────────────────────────────────────────────
\newcommand\ColumnTwo{%
\begin{panel}[slate]
\heading{5}{slate}{Filtering moves the centroids}
\includegraphics[width=\linewidth]{\fig{centroid_deviation_all_vitals_filter_mean}}
\caption{Centroid shift per vital after the all-vitals filter, by outcome
class (paper Fig.\ 3).}
\gap
Outliers are one-sided, so removing them moves the mean. Read across one
vital's bars: if the classes move by different amounts, the filter treated
them differently.
\end{panel}
\colgap
\begin{panel}[slate]
\heading{6}{slate}{\dots and it changes task difficulty}
\includegraphics{\fig{poster_task_divergence}}
\gap
Same features, folds and model; only the filter changes. Accuracy moves both
ways, often in opposite directions for the two tasks. The record-level filters
move it most, for a reason shown in \S10.
\end{panel}
\colgap
\begin{panel}[slate]
\heading{7}{slate}{Did the predictions change?}
\includegraphics{\fig{poster_mcnemar}}
\gap
McNemar's paired test ignores agreements and asks whether the filter fixed and
broke equally many patients. All arms use the baseline's threshold. No per-vital
filter is significant (all \textit{p}\,≥\,\PVitalMin).
\end{panel}
}

% ─────────────────────────────────────────────────────────────────────────────
% Column 3 — what the paper had no room for
% ─────────────────────────────────────────────────────────────────────────────
\newcommand\ColumnThree{%
\begin{panel}[rust]
\heading{8}{rust}{How many patients survive?}
\includegraphics{\fig{poster_survivors}}
\gap
The value-level filters edit \ReadingsEdited\ of \Readings\ readings
(\PctReadingsEdited) and remove no one. \emph{High invalid data} sounds lenient,
tolerating \HIDTolerated\ blank cells in 168, yet keeps \textbf{\NHID\ of \NRaw}.
\end{panel}
\colgap
\begin{panel}[rust]
\heading{9}{rust}{The survivors are not random}
\includegraphics{\fig{poster_prevalence}}
\gap
Record filters keep densely monitored stays, and ICU monitoring tracks acuity:
missingness is informative. So they select sicker patients. Prolonged stay goes
from a \PrevICURaw\ minority to a \textbf{\PrevICUHID\ majority}.
\end{panel}
\colgap
\begin{panel}[rust]
\heading{10}{rust}{The bars compare different people}
\includegraphics{\fig{poster_population}}
\gap
The \emph{unchanged} baseline already scores worse on each filter's survivors
(ICU, high invalid data: \BaseAccWholeICU\,→\,\BaseAccHIDicu). That is
\ShiftHIDicu\ of the published \DeltaAccHIDicu. On the same \NHID\ stays at one
threshold, the filter itself moves accuracy by \FixedDeltaHIDicu.
\end{panel}
\colgap
\begin{panel}[rust]
\heading{11}{rust}{McNemar only sees survivors}
\includegraphics{\fig{poster_mcnemar_sample}}
\gap
Deleted patients never enter the test. High invalid data shifts the ICU base
rate by \PrevShiftICUHID\ points, yet \p{\PHIDicu}: silence here is lost
power, not a clean bill of health.
\end{panel}
}

\begin{document}
\measurecolumn{1}{257}{\ColumnOne}
\measurecolumn{2}{257}{\ColumnTwo}
\measurecolumn{3}{257}{\ColumnThree}
\begin{tikzpicture}[remember picture, overlay]

% ── header ────────────────────────────────────────────────────────────────────
\fill[teal] ($(current page.north west)+(20mm,-22mm)$) rectangle ++(90mm,-3mm);
\node[anchor=north west, inner sep=0pt] at ($(current page.north west)+(20mm,-31mm)$)
  {\begin{minipage}[t]{600mm}
   {\serif\fontsize{104}{108}\selectfont\color{ink}Beyond the Model}\par\vspace{4mm}
   {\serif\fontsize{48}{56}\selectfont\color{bodytext}The critical role of data
    filtering in clinical machine learning}\par\vspace{8mm}
   {\fontsize{31}{38}\selectfont\color{ink}Noah Subedar*, Colin Campbell*, Wenjing
    Zhang, Dan Perri, Sarah Culgin, Andrew Hamilton-Wright}\par\vspace{3mm}
   {\fontsize{24}{30}\selectfont\color{muted}University of Guelph \textperiodcentered\
    St.\ Joseph's Healthcare Hamilton \textperiodcentered\ The Research Institute of
    St.\ Joe's Hamilton \textperiodcentered\ McMaster University
    \quad * equal contribution}
   \end{minipage}};

% Logos, if supplied: stacked in the header's right-hand corner.
\node[anchor=north east, inner sep=0pt] at ($(current page.north east)+(-20mm,-26mm)$)
  {\begin{minipage}[t]{170mm}\raggedleft
   \def\logo#1{\IfFileExists{Experiments/poster/logos/#1.pdf}%
     {\includegraphics[height=30mm,width=170mm,keepaspectratio]{Experiments/poster/logos/#1.pdf}\par\vspace{7mm}}%
     {\typeout{Missing logo: Experiments/poster/logos/#1.pdf}}}%
   \logo{guelph}\logo{stjoes}\logo{mcmaster}%
   \end{minipage}};

% ── hook: the finding, readable from across the room ─────────────────────────
\node[anchor=north west, inner sep=0pt] at ($(current page.north west)+(20mm,-136mm)$)
  {\begin{tcolorbox}[enhanced, width=801mm, height=82mm, colback=ink, frame hidden,
     arc=0pt, left=14mm, right=14mm, top=10mm, bottom=10mm, valign=center]
   \begin{minipage}[c]{252mm}
     \raggedright{\serif\fontsize{40}{48}\selectfont\color{white}Filtering is a
      modelling choice. Count what you removed.}
   \end{minipage}\hfill
   \newcommand\tile[3]{\begin{minipage}[c]{155mm}
     {\ttfamily\bfseries\fontsize{60}{64}\selectfont\color{#1}#2}\par\vspace{3mm}
     \raggedright{\fontsize{23}{28}\selectfont\color{white}#3}\end{minipage}}%
   \tile{teal!45}{0}{patients removed by the seven value-level filters}\hspace{12mm}%
   \tile{rust!55}{\RemovedHID}{of \NRaw\ stays removed by one record-level filter}\hspace{12mm}%
   \tile{rust!55}{\PrevICURaw\,→\,\PrevICUHID}{prolonged-stay prevalence
     among its survivors}
   \end{tcolorbox}};

% ── three columns ─────────────────────────────────────────────────────────────
\placecolumn{20}{257}{\ColumnOne}
\placecolumn{292}{257}{\ColumnTwo}
\placecolumn{564}{257}{\ColumnThree}

% ── bottom band ───────────────────────────────────────────────────────────────
\node[anchor=north west, inner sep=0pt] at ($(current page.north west)+(20mm,-989mm)$)
  {\begin{tcolorbox}[enhanced, width=529mm, height=128mm, colback=white, frame hidden,
     arc=0pt, boxrule=0pt, boxsep=0pt, left=10mm, right=10mm, top=9mm, bottom=8mm,
     borderline north={3mm}{0pt}{ink}, before upper={\body\raggedright}]
   \heading{12}{ink}{Three questions, asked together}
   {\small\renewcommand\arraystretch{1.3}
   \begin{tabularx}{\linewidth}{@{}>{\raggedright\arraybackslash\hsize=.62\hsize}X>{\raggedright\arraybackslash\hsize=1.06\hsize}X>{\raggedright\arraybackslash\hsize=1.16\hsize}X>{\raggedright\arraybackslash\hsize=1.16\hsize}X@{}}
   & \textbf{Centroid shift} & \textbf{McNemar vs.\ raw} & \textbf{Record accounting} \\[1mm]
   \level{teal}{all vitals} & moves; the classes move largely together
     & \p{\PAllVitalsMort} (mortality), \PAllVitalsICU\ (ICU): predictions unchanged
     & \leavevmode{\color{teal}0 stays removed; \ReadingsEdited\ readings edited} \\[2mm]
   \level{rust}{high invalid data} & large, and opposite in sign between the
     mortality classes
     & \p{\PHIDmort} (mortality), \PHIDicu\ (ICU), on only \NHID\ stays
     & \leavevmode{\color{rust}\RemovedHID\ stays removed; ICU prevalence \PrevICURaw\,→\,\PrevICUHID} \\
   \end{tabularx}}
   \gap
   Row one is artifact removal; row two is a different problem. The paper
   published the first two columns. The third costs one call to
   \texttt{len()}, and it explains the other two.
   \end{tcolorbox}};

\node[anchor=north west, inner sep=0pt] at ($(current page.north west)+(564mm,-989mm)$)
  {\begin{tcolorbox}[enhanced, width=257mm, height=128mm, colback=white, frame hidden,
     arc=0pt, boxrule=0pt, boxsep=0pt, left=10mm, right=10mm, top=9mm, bottom=8mm,
     borderline north={3mm}{0pt}{teal}, before upper={\body\raggedright}]
   \heading{13}{teal}{Take-aways}
   \newcommand\take[2]{\par\hangindent=11mm\hangafter=1%
     \makebox[11mm][l]{\ttfamily\bfseries\color{teal}#1}#2\par\vspace{5mm}}%
   \take{1}{\textbf{Filtering is part of the method,} not housekeeping.}
   \take{2}{\textbf{Level matters more than threshold.} \PctReadingsEdited\ of
     readings edited, no one removed, versus \PctRemovedHID\ of stays gone.}
   \take{3}{\textbf{Report the surviving n and its prevalence} beside every
     metric, and score the baseline on the same survivors.}
   \end{tcolorbox}};

% ── footer ────────────────────────────────────────────────────────────────────
\node[anchor=north west, inner sep=0pt] at ($(current page.north west)+(20mm,-1127mm)$)
  {\begin{minipage}[t]{801mm}
   \color{rule}\rule{801mm}{1.2mm}\par\vspace{4mm}
   \fontsize{17}{22}\selectfont\color{muted}%
   \begin{minipage}[t]{\ifx\RepoURL\empty 801mm\else 740mm\fi}
   \textbf{Scope.} MIMIC-III (\NRaw\ ICU stays), mean hourly aggregation,
   5-fold\,×\,4 cross-validation; McNemar at the baseline's decision threshold.
   One database, seven vitals, two tasks, one classifier.
   \textbf{References.} Johnson et al., MIMIC-III, \emph{Sci.\ Data} 2016 ·
   Wang et al., MIMIC-Extract, \emph{CHIL} 2020 · McNemar, \emph{Psychometrika} 1947 ·
   Youden, \emph{Cancer} 1950.
   \textbf{Code.} \ifx\RepoURL\empty available on GitHub\else\texttt{\RepoURL}\fi.
   Computation on nibi, Digital Research Alliance of Canada.
   Contact: \texttt{wzhang25@uoguelph.ca}
   \end{minipage}%
   \ifx\RepoURL\empty\else\hfill\qrcode[height=38mm]{\RepoURL}\fi
   \end{minipage}};

\end{tikzpicture}
\end{document}
